\begin{proof}[Proof of \cref{obj:lem:signed-depth-bootstrap-factor}]
Let
\[
        \alpha=e^{-1/t_0},\qquad c_0=\frac{1-\alpha}{4},\qquad
        \eta_0=\frac{\alpha}{1-c_0}.
\]
We prove the stated uniform bound by constructing the constant explicitly.

\begin{enumerate}
\item Since \(t_0>0\), one has \(0<\alpha<1\). Hence
\[
        0\le c_0=\frac{1-\alpha}{4}\le \frac14,
        \qquad
        1-c_0=\frac{3+\alpha}{4}>0,
\]
and therefore
\[
        0\le \eta_0=\frac{\alpha}{1-c_0}<1,
\]
because \(\alpha<1-c_0\) is equivalent to \(\alpha<(3+\alpha)/4\), or \(\alpha<1\). % lean: CausalSmith.Stat.PomdpLatentOverlapMinimax.exists_uniform_signedDepth_bootstrap_factor

\item The sequence \(Q\eta_0^Q\) is bounded above. To see this in the same normalization, set
\[
        \bar\eta=\frac{1+\eta_0}{2}.
\]
Then \(0\le \bar\eta<1\) and \(|\eta_0|<\bar\eta\). Writing
\[
        Q\eta_0^Q=\bigl(Q(\eta_0/\bar\eta)^Q\bigr)\bar\eta^Q,
\]
the first factor tends to \(0\) by the ratio test, since
\[
        \frac{(Q+1)(\eta_0/\bar\eta)^{Q+1}}
             {Q(\eta_0/\bar\eta)^Q}
        =\Bigl(1+\frac1Q\Bigr)\frac{\eta_0}{\bar\eta}
        \longrightarrow \frac{\eta_0}{\bar\eta}<1,
\]
while \(\bar\eta^Q\to0\). Thus \(Q\eta_0^Q\to0\), and in particular there is a finite \(M_\star\ge0\) such that
\[
        Q\eta_0^Q\le M_\star\qquad(Q\in\mathbb N).
\]
 % lean: CausalSmith.Stat.PomdpLatentOverlapMinimax.exists_uniform_signedDepth_bootstrap_factor

\item Define
\[
        K_0=\exp(2c_0M_\star).
\]
Then \(K_0>0\). Fix \(Q\in\mathbb N\) and put
\[
        x_Q=c_0\eta_0^Q.
\]
From \(0\le c_0\le1/4\) and \(0\le\eta_0<1\), we have
\[
        0\le x_Q\le \frac14,
        \qquad
        1-x_Q>0.
\]
For every \(x\in[0,1/4]\),
\[
        (1-x)^{-1}\le 1+2x\le e^{2x}.
\]
Indeed, the first inequality follows after multiplication by \(1-x>0\), since
\[
        (1+2x)(1-x)=1+x-2x^2\ge1,
\]
and the second is the standard bound \(1+y\le e^y\) with \(y=2x\). Applying this with \(x=x_Q\) and raising to the \(Q\)-th power gives
\[
        \bigl((1-x_Q)^{-1}\bigr)^Q
        \le \bigl(e^{2x_Q}\bigr)^Q
        =e^{2Qx_Q}.
\]
 % lean: CausalSmith.Stat.PomdpLatentOverlapMinimax.exists_uniform_signedDepth_bootstrap_factor

\item Finally,
\[
        2Qx_Q=2c_0\,Q\eta_0^Q\le 2c_0M_\star,
\]
so monotonicity of the exponential yields
\[
        e^{2Qx_Q}\le e^{2c_0M_\star}=K_0.
\]
Combining the preceding displays,
\[
\frac{1}{\bigl(1-c_0\eta_0^Q\bigr)^Q}
=
\bigl((1-x_Q)^{-1}\bigr)^Q
\le
K_0.
\]
Since \(Q\in\mathbb N\) was arbitrary, this \(K_0\) proves the claim. % lean: CausalSmith.Stat.PomdpLatentOverlapMinimax.exists_uniform_signedDepth_bootstrap_factor
\end{enumerate}
\end{proof}
